\renewcommand{\Anno}{1999}

\renewcommand{\NumeroProblema}{2}

\renewcommand{\Testo}{Consideriamo un lago circolare di 1 km di
diametro.  Sia $AB$ un diametro e sia $C$ il punto medio di una delle semicirconferenze delimitate da $A$ e $B$.  Un istituto di ricerca vuole costruire una piattaforma sul lago, posizionata in un punto del diametro $AB$.  Per rendere operativa la piattaforma occorrono due tipi di collegamento, entrambi realizzati dalla stessa ditta:
\begin{itemize}
	\item cavi elettrici, collegati in linea retta con $C$ (costo: 1 euro al metro);
	
	\item cavi con fibre ottiche, collegati in linea retta con il punto più vicino tra $A$ e $B$ (costo: 5 euro al metro).
\end{itemize}
	
Uno sponsor si è offerto di pagare la minore tra le due spese di cui ai punti precedenti; la rimanente spesa graverà sulle casse dell'istituto.
	
\begin{enumerate}
	\item A quale distanza dal centro del lago conviene che l'istituto di ricerca posizioni la piattaforma, se vuole minimizzare le spese a suo carico?
	
	\item Quale posizione sarebbe la più conveniente per la ditta che posiziona i cavi (cioè le farebbe guadagnare di più)?

\end{enumerate}}

\renewcommand{\Risposta}{\begin{enumerate}
	\item  375 metri.
	
	\item  Il centro del lago.
\end{enumerate}}

\renewcommand{\Aiutino}{Detta $x$ la distanza dal centro, si
determinano in funzione di $x$ il costo dei cavi elettrici e delle fibre ottiche.  A questo punto la spesa a carico dell'istituto è il massimo tra i due costi: per il punto (a) non resta dunque che tracciarne il grafico (si tratta del grafico del massimo tra due funzioni) e vedere per quale valore di $x$ assume il minimo.  Il guadagno della ditta è invece la somma dei due costi.  Per il punto (b) occorre quindi trovare il massimo della somma delle due funzioni: per farlo si può ricorrere al calcolo differenziale, oppure indovinare ad occhio il punto di massimo e dimostrarlo
risolvendo una disequazione.}

\renewcommand{\SoluzioneUfficiale}{
}

\renewcommand{\SoluzioniAlternative}{
}

\renewcommand{\Osservazioni}{
}

\renewcommand{\Commenti}{
}
% Aggiunta redazionale: la soluzione preesistente non viene modificata.
\providecommand{\SoluzioneRedazionale}{}
\renewcommand{\SoluzioneRedazionale}{
% SSSUP-BOZZA-AI: SoluzioneRedazionale; autore=ChatGPT; data=2026-11-08; stato=DA-VERIFICARE
% SSSUP-PROVENIENZA: AI: elaborazione nuova dal testo
Misuriamo le lunghezze in metri e indichiamo con $x\in[0,500]$ la distanza della piattaforma dal centro. Per simmetria non importa su quale metà del diametro si trovi. I due costi sono
\[
E(x)=\sqrt{500^2+x^2},\qquad F(x)=5(500-x).
\]

\paragraph{Punto (a).}
Lo sponsor paga il minore dei due costi, quindi l'istituto paga $\max\{E(x),F(x)\}$. La funzione $E$ è crescente, $F$ è decrescente, ed esse si incontrano in un unico punto. Il massimo ha il suo minimo proprio in tale punto. L'equazione
\[
500^2+x^2=25(500-x)^2
\]
ha radici $375$ e $2000/3$; solo $375$ appartiene all'intervallo ammesso. Il costo minimo dell'istituto è allora $625$ euro.

\paragraph{Punto (b).}
La ditta riceve la somma $E(x)+F(x)$, indipendentemente da chi paga. Per $x\geq0$ si ha $\sqrt{500^2+x^2}\leq500+x$, quindi
\[
E(x)+F(x)\leq3000-4x\leq3000.
\]
Il massimo $3000$ euro si ottiene solo per $x=0$: la piattaforma deve stare nel centro.
}
